\chapter*{图片致谢}

本书中的图示均为原创。照片依其各自的自由许可使用，谨向其作者致谢：

\begin{itemize}
  \item 第~\ref{ch:b3:genomics} 章：弗雷德里克·桑格的照片——
        公有领域（维基共享资源）。
  \item 第~\ref{ch:b3:genetic-engineering} 章：黄金大米与普通大米并列，国际水稻研究所，2011 年——
        CC~BY~2.0（维基共享资源）。
  \item 第~\ref{ch:b3:structural-biology} 章：马克斯·佩鲁茨的照片，美联社，1962 年——公有领域（维基共享资源）。
  \item 第~\ref{ch:b3:bacteriology} 章：罗伯特·科赫的肖像，Wellcome Collection——
        CC~BY~4.0（维基共享资源）；亚历山大·弗莱明在其实验室中的照片，美国海军医学与外科局档案——
        无已知版权限制（维基共享资源）。
  \item 第~\ref{ch:b3:innate-immunity} 章：埃利·梅契尼科夫的照片，Bain News Service，美国国会图书馆——公有领域（维基共享资源）。
  \item 第~\ref{ch:b3:adaptive-immunity} 章：爱德华·詹纳的雕刻肖像，1807 年——公有领域（维基共享资源）。
  \item 第~\ref{ch:b3:nervous-systems} 章：Santiago Ramón y Cajal 绘制的浦肯野细胞图，1899 年——公有领域（维基共享资源）。
  \item 第~\ref{ch:b3:endocrinology} 章：弗雷德里克·班廷与查尔斯·贝斯特的照片，约 1924 年——公有领域（维基共享资源）。
  \item 第~\ref{ch:b3:molecular-evolution} 章：木村资生的照片，1986 年，摄影者不详，经 S.~Kumar 与 T.~Gojobori，\emph{Molecular Biology and Evolution}（2023）转引——CC~BY~4.0（维基共享资源）。
  \item 第~\ref{ch:b3:behavioural-ecology} 章：尼科·廷伯根的照片，Rob Mieremet 摄，Anefo，1973 年——CC0，Nationaal Archief（维基共享资源）。
  \item 第~\ref{ch:b3:conservation-biology} 章：最后一只旅鸽玛莎的照片，Enno Meyer，1912 年——公有领域（维基共享资源）；黄石国家公园的狼，美国国家公园管理局，1996 年——公有领域（维基共享资源）；卡卡坡鹦鹉 Sirocco 的照片，新西兰自然保护部，2012 年——CC~BY~2.0（维基共享资源）。
\end{itemize}

每张照片的完整来源链接和许可证网址，均列于本书代码仓库中的 \texttt{images/book5/CREDITS.md}。

\bigskip
与照片及原创图示并列，全书各章中约九十幅插图为 AI 生成图像，由 OpenAI 图像生成功能依本书编辑撰写的提示词生成，并在纳入前经过生物学准确性审校。每幅生成图像的完整提示词均记录于代码仓库中的 \texttt{images/book5/ai/PROMPTS.md}。
\clearpage
